\documentclass[12pt,a4paper]{article}
\usepackage[utf8]{inputenc}
\usepackage[indonesian]{babel}
\usepackage{amsmath}
\usepackage{amsfonts}
\usepackage{amssymb}
\usepackage{geometry}

\geometry{margin=2.5cm}

\title{Menghitung Luas Permukaan Benda Putar Sumbu-y}
\author{Ahmad Fakhrul Bawani}
\date{}

\begin{document}

\maketitle

\section*{Soal}

Find the area of the surface generated by revolving the curve:
\begin{equation*}
  x = \frac{y^3}{5}, \quad \text{untuk } 0 \le y \le 5, \quad \text{mengelilingi sumbu-}y
\end{equation

\subsection*{1. Rumus Umum}
Rumus luas permukaan ($S$) benda putar yang dihasilkan oleh perputaran mengelilingi sumbu-$y$ untuk fungsi dalam variabel $y$ dari $y = c$ sampai $y = d$ didefinisikan sebagai:
\begin{equation*}
  S = 2\pi \int_{c}^{d} x \cdot \sqrt{1 + \left(\frac{dx}{dy}\right)^2} \, dy
\end{equation*}

\subsection*{2. Menentukan Turunan Fungsi}
Cari turunan pertama fungsi $x = \frac{1}{5}y^3$ terhadap $y$:
\begin{equation*}
  \frac{dx}{dy} = \frac{1}{5}(3y^2) = \frac{3y^2}{5}
\end{equation*}

Kuadratkan hasil turunan tersebut:
\begin{equation*}
  \left(\frac{dx}{dy}\right)^2 = \left(\frac{3y^2}{5}\right)^2 = \frac{9y^4}{25}
\end{equation*}

Masukkan ke dalam komponen akar:
\begin{equation*}
  \sqrt{1 + \left(\frac{dx}{dy}\right)^2} = \sqrt{1 + \frac{9y^4}{25}} = \sqrt{\frac{25 + 9y^4}{25}} = \frac{1}{5}\sqrt{25 + 9y^4}
\end{equation*}

\subsection*{3. Solusi Integrasi}
Substitusikan komponen $x = \frac{y^3}{5}$ dan bentuk akar di atas ke dalam rumus luas permukaan dengan batas $c = 0$ dan $d = 5$:
\begin{align*}
  S &= 2\pi \int_{0}^{5} \left(\frac{y^3}{5}\right) \cdot \left(\frac{1}{5}\sqrt{25 + 9y^4}\right) \, dy \\
  S &= \frac{2\pi}{25} \int_{0}^{5} y^3 (25 + 9y^4)^{\frac{1}{2}} \, dy
\end{align*}

Gunakan metode substitusi untuk menyelesaikan integral tersebut:
\begin{align*}
  \text{Misalkan: } & u = 25 + 9y^4 \implies du = 36y^3 \, dy \implies y^3 \, dy = \frac{1}{36} \, du \\
  \text{Batas baru: } & \text{Untuk } y = 0 \implies u = 25 + 9(0)^4 = 25 \\
  & \text{Untuk } y = 5 \implies u = 25 + 9(5)^4 = 25 + 9(625) = 5650
\end{align*}

Substitusikan variabel baru $u$ ke dalam integral:
\begin{align*}
  S &= \frac{2\pi}{25} \int_{25}^{5650} u^{\frac{1}{2}} \cdot \left(\frac{1}{36} \, du\right) \\
  S &= \frac{2\pi}{900} \int_{25}^{5650} u^{\frac{1}{2}} \, du \\
  S &= \frac{\pi}{450} \left[ \frac{2}{3}u^{\frac{3}{2}} \right]_{25}^{5650} \\
  S &= \frac{\pi}{675} \left[ u\sqrt{u} \right]_{25}^{5650} \\
  S &= \frac{\pi}{675} \left( 5650\sqrt{5650} - 25\sqrt{25} \right)
\end{align*}

Sederhanakan bentuk akar $\sqrt{5650} = \sqrt{25 \cdot 226} = 5\sqrt{226}$:
\begin{align*}
  5650\sqrt{5650} &= 5650 \cdot 5\sqrt{226} = 28250\sqrt{226} \\
  25\sqrt{25} &= 25 \cdot 5 = 125
\end{align*}

Substitusikan kembali nilai numerik tersebut ke dalam persamaan:
\begin{align*}
  S &= \frac{\pi}{675} \left( 28250\sqrt{226} - 125 \right) \\
  S &= \frac{125\pi}{675} \left( 226\sqrt{226} - 1 \right)
\end{align*}

Sederhanakan pecahan $\frac{125}{675}$ dengan membagi pembilang dan penyebut dengan $25$:
\begin{equation*}
  S = \frac{5}{27}\pi \left( 226\sqrt{226} - 1 \right)
\end{equation*}

Jadi, luas permukaan benda putar tersebut adalah \textbf{$\dfrac{5}{27}\pi \left( 226\sqrt{226} - 1 \right)$ satuan luas}.

\end{document}